\subsection{Four familiar shapes, one geometric problem}

The four curves differ in their Cartesian equations, but the decisive geometry
is already visible:

\[
\begin{array}{c|c|c}
\text{curve} & \text{distinguished data} & \text{large-scale behaviour}\\
\hline
\text{circle} & \text{one centre and one radius} & \text{closed}\\
\text{ellipse} & \text{two foci, fixed sum} & \text{closed}\\
\text{parabola} & \text{one focus and one directrix} & \text{one open branch}\\
\text{hyperbola} & \text{two foci, fixed difference} & \text{two open branches}
\end{array}
\]

The centre is the natural origin for a circle. It is also convenient for a
Cartesian ellipse or hyperbola. But it is not the natural point from which to
study focal distance. For that purpose one should move the origin to a focus.

\begin{corebox}[The need for a new coordinate system]
From a focus, a point on a conic is determined by two questions:
\begin{enumerate}
    \item In which direction do we look?
    \item How far must we travel in that direction before reaching the curve?
\end{enumerate}
These are exactly the questions answered by polar coordinates.
\end{corebox}

\begin{center}
\begin{tikzpicture}[scale=0.9,>=stealth]
    \draw[axis,->] (-4.8,0)--(5.1,0) node[right] {$x$};
    \draw[axis,->] (0,-3.0)--(0,3.2) node[above] {$y$};
    \draw[subset] (0.6,0) ellipse (3.8 and 2.1);
    \coordinate (F) at (-2.55,0);
    \coordinate (P1) at (2.4,1.55);
    \coordinate (P2) at (-0.4,2.02);
    \coordinate (P3) at (-3.0,-0.95);
    \draw[blue!70!black,line width=1pt,->] (F)--(P1);
    \draw[blue!70!black,line width=1pt,->] (F)--(P2);
    \draw[blue!70!black,line width=1pt,->] (F)--(P3);
    \fill[geometricSubsetColor] (F) circle (2.3pt) node[below] {chosen focus};
    \node[align=left] at (3.8,-2.35)
        {different directions\\different focal distances};
\end{tikzpicture}
\end{center}

The next section constructs this coordinate language before using it on conics.
